% !TEX root = main.tex
% 《杂阿含经·仿罗什风新译》研究译注本 — 导言区

\usepackage[svgnames]{xcolor}
\usepackage[bookmarksnumbered=true,unicode=true]{hyperref}
\hypersetup{
  pdftitle={杂阿含经·仿罗什风新译·研究译注本（大正卷序）},
  pdfauthor={Agama–Kumarajiva Project},
  colorlinks=true,
  linkcolor=MidnightBlue,
  citecolor=MidnightBlue,
  urlcolor=MidnightBlue
}

\usepackage{geometry}
\geometry{a4paper, margin=2.3cm, headheight=15pt}

\usepackage{setspace}
\setstretch{1.28}

\usepackage{enumitem}
\setlist{nosep}

\usepackage{graphicx}
\definecolor{coverbg}{HTML}{C9B896}
\definecolor{coverink}{HTML}{5C4A32}

\usepackage{longtable}
\usepackage{array}
\usepackage{booktabs}
\usepackage{fancyhdr}
\usepackage{titlesec}

% Fandol 缺大量繁体字；fontset=none 后显式指定 Noto CJK
\setCJKmainfont[
  Path = /usr/share/fonts/google-noto-serif-cjk-fonts/,
  Extension = .ttc,
  UprightFont = NotoSerifCJK-Regular,
  BoldFont = NotoSerifCJK-Bold,
]{NotoSerifCJK}
\setCJKsansfont[
  Path = /usr/share/fonts/google-noto-sans-cjk-fonts/,
  Extension = .ttc,
  UprightFont = NotoSansCJK-Regular,
  BoldFont = NotoSansCJK-Bold,
]{NotoSansCJK}
\setmainfont{Liberation Serif}
\setsansfont{Liberation Sans}
\setmonofont{Liberation Mono}
% 省略号、重建标记等易被当作西文缺字；归入 CJK 字类
\xeCJKDeclareCharClass{CJK}{"22EF,"25C7}

% ---------- 研究本专用环境 ----------
% 卷：不用 ctexbook 自带「第X章」编号
\newcommand{\juan}[1]{%
  \clearpage
  \chapter*{#1}%
  \markboth{#1}{#1}%
  \addcontentsline{toc}{chapter}{#1}%
  \vspace{0.5em}%
}

% 一卷内多个相应时分节
\newcommand{\samyukta}[1]{%
  \section*{#1}%
  \addcontentsline{toc}{section}{#1}%
  \vspace{0.3em}%
}

% 经题：整卷单一相应时用 section（避免 subsection 书签错乱嵌套）
\newcommand{\sattitlesec}[1]{%
  \section*{#1}%
  \addcontentsline{toc}{section}{#1}%
  \vspace{0.3em}%
}

\newcommand{\sattitle}[1]{%
  \subsection*{#1}%
  \addcontentsline{toc}{subsection}{#1}%
  \vspace{0.3em}%
}

\newcommand{\suttahrule}{\par\vspace{0.6em}\hrule\vspace{1.2em}}

% 页眉
\pagestyle{fancy}
\fancyhf{}
\fancyhead[LE,RO]{\small\itshape 杂阿含经·仿罗什风新译}
\fancyhead[LO,RE]{\small\leftmark}
\fancyfoot[C]{\thepage}
\renewcommand{\headrulewidth}{0.4pt}

% 扉页：无页眉页脚
\fancypagestyle{coverpage}{%
  \fancyhf{}%
  \renewcommand{\headrulewidth}{0pt}%
  \renewcommand{\footrulewidth}{0pt}%
}

\titleformat{\chapter}[display]
  {\normalfont\huge\bfseries}{\chaptertitlename\ \thechapter}{18pt}{\Huge}
